\uuid{5wy2}
\titre{Convergence}
\auteur{bodin ; hanani ; mzari}
\organisation{exo7}
\licence{CC BY-NC-SA 4.0}
\datecreate{2020-01-20}
\niveau{L1}
\module{Analyse}
\chapitre{Suite}
\sousChapitre{Autre}
\difficulte{2}
\theme{Suites, convergence}
\exo7qcmid{lille-1.suites.019}

\contenu{
\begin{question}
\qclassification{121.00}
\qtags{motcle=convergence}
Soit $(u_n)$ une suite réelle. On suppose que $\displaystyle u_n\geq \sqrt{n}$ pour tout $n\geq 0$. Que peut-on en déduire ?
\begin{answers}
    \good{La suite $(u_n)$ n'est pas majorée.}
    \bad{La suite $(u_n)$ est croissante.}
    \bad{La suite $(u_n)$ est convergente.}
    \good{$\displaystyle \lim _{n\to +\infty}u_n=+\infty$}
\end{answers}
\begin{explanations}
Si $(u_n)$ était majorée, il en serait de même pour $\sqrt{n}$ ce qui est absurde. Donc $(u_n)$ est une suite non majorée. Par passage à la limite, on a :
$$\lim _{n\to +\infty}u_n\geq \lim _{n\to +\infty}\sqrt{n}=+\infty\Rightarrow \displaystyle \lim _{n\to +\infty}u_n=+\infty.$$
\end{explanations}
\end{question}
}
